\hypertarget{appendix-343615f49ddf}{}
\section{IEEE-754 floating point}\label{appendix-343615f49ddf-ieee-754-floating-point}

\textbf{Curriculum depth:} Developed · \textbf{First seen:} C001 · \textbf{Latest development:} C005

\subsection{Definition and mechanism}\label{appendix-343615f49ddf-definition-and-mechanism}

IEEE-754 specifies floating formats and behavior including representation classes, rounding directions and exception conditions. A finite bit pattern is interpreted through sign, exponent and significand rules, not as a two's-complement integer.

Binary32 has 1/8/23 stored fields and 24-bit normal precision; binary64 has 1/11/52 fields and 53-bit normal precision. Word 412A0000 encodes 10.625 in binary32. Exponent-zero and all-one fields distinguish zero/subnormal and infinity/NaN. Normal spacing expands with exponent; nearest-even and directed modes choose representable results under explicit rules. These format names do not prove every C++ float/double has those representations.

M005 develops alignment, cancellation, FMA and accumulated error. Numeric equality is not representation identity, and moving rounding points can change the computed result.

\subsection{Prerequisites}\label{appendix-343615f49ddf-prerequisites}

\begin{itemize}
\tightlist
\item
  \hyperref[appendix-ac697f5c66e8]{Bit field} --- A field is a declared range with offset s and width w, not automatically a C++ language bit-field member.
\item
  \hyperref[appendix-f414bd4ed05c]{Interpretation contract} --- For \texttt{0xAD}, unsigned interpretation gives 173; eight-bit two's complement gives −83.
\end{itemize}

\subsection{Used by}\label{appendix-343615f49ddf-used-by}

\begin{itemize}
\item
  \hyperref[appendix-34443cec7170]{Binary16} --- uses this mechanism to define its representation, error or validation contract.
\item
  \hyperref[appendix-18b8dd966b5d]{Bfloat16} --- uses this mechanism to define its representation, error or validation contract.
\item
  \hyperref[appendix-41e8ef22a21e]{Fused multiply-add} --- Fused operation semantics and destination rounding are part of the arithmetic contract.
\item
  \hyperref[appendix-5057f9ad327f]{Binary32} --- uses ieee-754 floating point to establish this mechanism's representation or policy.
\item
  \hyperref[appendix-ed8531b46754]{Floating-point exception flags} --- uses ieee-754 floating point to establish this mechanism's representation or policy.
\item
  \hyperref[appendix-34ee45f161c1]{Bit casting and representation} --- uses ieee-754 floating point to establish this mechanism's representation or policy.
\item
  \hyperref[appendix-7f1e17dcafb1]{Robust norm and intermediate range} --- uses ieee-754 floating point to establish this mechanism's representation or policy.
\end{itemize}

Connect later dependent concepts here when their lessons arrive.

\subsection{Recall and next depth}\label{appendix-343615f49ddf-recall-and-next-depth}

Reconstruct 10.625's fields, distinguish epsilon from minimum positive values, and resolve an exact midpoint without calling it decimal rounding.

\textbf{Pending:} Extend this mechanism to later algorithms and target-specific behavior. See the M005 chapter for its current examples and validity assumptions.

\subsection{Source and navigation}\label{appendix-343615f49ddf-source-and-navigation}

\hyperref[01-foundations-i-ch001-when-arithmetic-meets-uncertainty]{M005 development}

\hyperref[01-foundations-i-ch001-a-stretching-ruler-for-real-valued-quantities]{C004 development}

\hyperref[01-foundations-i-ch001-from-a-signal-to-an-interpretation]{C001 teaching chapter} · \hyperref[appendix-fd090d982a15]{Part definitions and mechanisms} · \href{../02-Foundations-II/main.pdf\#nameddest=appendix-4ad47bdd269a}{Dependency Map (02-Foundations-II, appendix-4ad47bdd269a)}
